%2multibyte Version: 5.50.0.2953 CodePage: 1252

\documentclass[notes=show]{beamer}
%%%%%%%%%%%%%%%%%%%%%%%%%%%%%%%%%%%%%%%%%%%%%%%%%%%%%%%%%%%%%%%%%%%%%%%%%%%%%%%%%%%%%%%%%%%%%%%%%%%%%%%%%%%%%%%%%%%%%%%%%%%%%%%%%%%%%%%%%%%%%%%%%%%%%%%%%%%%%%%%%%%%%%%%%%%%%%%%%%%%%%%%%%%%%%%%%%%%%%%%%%%%%%%%%%%%%%%%%%%%%%%%%%%%%%%%%%%%%%%%%%%%%%%%%%%%
\usepackage{amsfonts}
\usepackage{graphicx}
\usepackage{amssymb}
\usepackage{amsmath}
\usepackage{mathpazo}
\usepackage{hyperref}
\usepackage{multimedia}

\setcounter{MaxMatrixCols}{10}
%TCIDATA{OutputFilter=LATEX.DLL}
%TCIDATA{Version=5.50.0.2953}
%TCIDATA{Codepage=1252}
%TCIDATA{<META NAME="SaveForMode" CONTENT="2">}
%TCIDATA{BibliographyScheme=Manual}
%TCIDATA{Created=Wednesday, May 01, 2013 14:54:39}
%TCIDATA{LastRevised=Wednesday, September 28, 2022 11:16:34}
%TCIDATA{<META NAME="GraphicsSave" CONTENT="32">}
%TCIDATA{<META NAME="DocumentShell" CONTENT="Other Documents\SW\Slides - Beamer">}
%TCIDATA{Language=American English}
%TCIDATA{CSTFile=beamer.cst}
%TCIDATA{PageSetup=72,72,72,83,0}
%TCIDATA{Counters=arabic,1}
%TCIDATA{AllPages=
%H=36
%F=36
%}


\newenvironment{stepenumerate}{\begin{enumerate}[<+->]}{\end{enumerate}}
\newenvironment{stepitemize}{\begin{itemize}[<+->]}{\end{itemize} }
\newenvironment{stepenumeratewithalert}{\begin{enumerate}[<+-| alert@+>]}{\end{enumerate}}
\newenvironment{stepitemizewithalert}{\begin{itemize}[<+-| alert@+>]}{\end{itemize} }
\input{tcilatex}
\usetheme{JuanLesPins}
\usecolortheme{whale}

\begin{document}

\title{Econometrics Lecture 1: \\
Introduction to Econometrics}
\author{Katerina Petrova \ \ \ \ \ \ \ \ \ \ \ \ \ \ \ \ \ \ \ \ \ \ \ \ \ \
\ \ \ \ \ \ \ \ \ \ \ \ \ \ \ \ \ \ \ \ \ \ \ \ \ \ \ \ \ \ \ \ \ \ \ \ \ \
\ \ \ \ \ \ \ \ \ \ \ \ \ \ \ \ \ \ \ \ \ \ \ }
\date{}
\maketitle

\section{About the Course}

%TCIMACRO{\TeXButton{BeginFrame}{\begin{frame}}}%
%BeginExpansion
\begin{frame}%
%EndExpansion

\QTR{frametitle}{Preliminaries}

\begin{itemize}
\item Dr. Katerina Petrova

\item Office: 20.1E74

\item E-mail: Katerina.Petrova@upf.edu

\item Tutorials/Labs: taught by Valeria Gargiulo

\item Email: valeria.gargiulo@upf.edu

\item Office hour: Wednesdays 14:00-15:00

\item Course material: Aula Global
\end{itemize}

%TCIMACRO{\TeXButton{EndFrame}{\end{frame}}}%
%BeginExpansion
\end{frame}%
%EndExpansion

%TCIMACRO{\TeXButton{BeginFrame}{\begin{frame}}}%
%BeginExpansion
\begin{frame}%
%EndExpansion

\QTR{frametitle}{Topics to cover}

\begin{enumerate}
\item The classical linear regression model (3-4 weeks)

\begin{itemize}
\item OLS and the Gauss Markov Theorem

\item Frisch-Waugh-Lovell theorem

\item asymptotic results/ hypothesis testing.
\end{itemize}

\item Endogeneity and Instrumental Variables estimator (1 week)

\item Heteroskedasticity and GLS (1 week)

\item Nonlinear Estimators (Maximum Likelihood) (2 weeks)

\item Introduction to Time Series (1 weeks)\ if time permits

\begin{itemize}
\item Covariance stationary processes, ARMA models

\item Multivariate extensions (VAR models)
\end{itemize}
\end{enumerate}

%TCIMACRO{\TeXButton{EndFrame}{\end{frame}}}%
%BeginExpansion
\end{frame}%
%EndExpansion

%TCIMACRO{\TeXButton{BeginFrame}{\begin{frame}}}%
%BeginExpansion
\begin{frame}%
%EndExpansion

\QTR{frametitle}{Reading:}

\begin{itemize}
\item Lecture Notes

\item Main Textbook: Bruce Hansen \textit{Econometrics} 2019

\item Open access resource and available online at:

https://perhuaman.files.wordpress.com/2014/06/econometrics-bruce-hansen-2014.pdf
\end{itemize}

%TCIMACRO{\TeXButton{EndFrame}{\end{frame}}}%
%BeginExpansion
\end{frame}%
%EndExpansion

%TCIMACRO{\TeXButton{BeginFrame}{\begin{frame}}}%
%BeginExpansion
\begin{frame}%
%EndExpansion

\QTR{frametitle}{Maths/Stats}

\begin{itemize}
\item Elementary level of matrix algebra and statistics is assumed

\item For a revision, please refer to 

\begin{itemize}
\item The two sets of lecture slides on Aula Global: Matrix Algebra Revision
and Probability Revision 

\item Appendix A and B in the textbook
\end{itemize}
\end{itemize}

%TCIMACRO{\TeXButton{EndFrame}{\end{frame}}}%
%BeginExpansion
\end{frame}%
%EndExpansion

%TCIMACRO{\TeXButton{BeginFrame}{\begin{frame}}}%
%BeginExpansion
\begin{frame}%
%EndExpansion

\QTR{frametitle}{Assessment}

\begin{itemize}
\item 10\% - Marked Problem Set, PS3 (allowed to work in groups of max 3)

\item 30\% - Class Test: Week 7 or Week 8 (9th November or 14/16 November)?

\item 60\% - Final Exam
\end{itemize}

%TCIMACRO{\TeXButton{EndFrame}{\end{frame}}}%
%BeginExpansion
\end{frame}%
%EndExpansion

\section{What is Econometrics about?}

%TCIMACRO{\TeXButton{BeginFrame}{\begin{frame}}}%
%BeginExpansion
\begin{frame}%
%EndExpansion

\QTR{frametitle}{What is econometrics?}

\begin{itemize}
\item It's the study of the properties and applications of statistical tools
to economic or financial data

\item We need statistical tools as all economic/financial data are
observational not experimental

\item We treat series as random variables and adopt a probability approach

\item We assume that there is a true data generating process given, for
example, by a density $f_{X}\left( x;\theta \right) $

\item Even if we know this DGP, we only observe a sample of realisations $%
\left( X_{1},...,X_{n}\right) $ of $x$

\item Moreover, the DGP depends on an unknown parameter vector $\theta $
\end{itemize}

%TCIMACRO{\TeXButton{EndFrame}{\end{frame}}}%
%BeginExpansion
\end{frame}%
%EndExpansion

%TCIMACRO{\TeXButton{BeginFrame}{\begin{frame}}}%
%BeginExpansion
\begin{frame}%
%EndExpansion

\QTR{frametitle}{Point Estimation}

\begin{itemize}
\item The idea of point estimation is to use data in order to estimate $%
\theta $, i.e., choose a random variable $\hat{\theta}=g\left(
X_{1},...,X_{n}\right) $ based on a sample of $n$ observations from $f\left(
x;\theta \right) $ in order to make a \textquotedblleft good
guess\textquotedblright\ for $\theta $
\end{itemize}

\begin{definition}
Given an unknown parameter $\theta $, the \textbf{parameter space} $\Theta $
is the set containing all possible values of $\theta $. Typically, $\Theta
\subseteq \mathbb{R}^{k}$.
\end{definition}

\begin{definition}
A \textbf{statistic} based on a sample of observations $\left(
X_{1},...,X_{n}\right) $ is any function $g\left( X_{1},...,X_{n}\right) $
of the random variables $X_{1},...,X_{n}$.
\end{definition}

%TCIMACRO{\TeXButton{EndFrame}{\end{frame}}}%
%BeginExpansion
\end{frame}%
%EndExpansion

%TCIMACRO{\TeXButton{BeginFrame}{\begin{frame}}}%
%BeginExpansion
\begin{frame}%
%EndExpansion

\QTR{frametitle}{Point Estimation}

\begin{itemize}
\item If the observations come from a family of distributions $f\left(
X_{1},...,X_{n};\theta \right) $ indexed by an unknown parameter $\theta ,$
an \textbf{estimator} $\hat{\theta}$\ of $\theta $ is a statistic used to
approximate $\theta $.

\item If $\theta $ is a vector of unknown parameters, $\hat{\theta}$ is a
vector of estimators. We sometimes write an estimator as $\hat{\theta}_{n}$
in order to emphasize its dependence on the sample size $n$.

\item Note that, by definition, an estimator $\hat{\theta}_{n}$ is a random
variable (or vector) that depends on the random variables $X_{1},...,X_{n}$
only: $\hat{\theta}_{n}$ \textit{does not depend on} $\theta $.
\end{itemize}

%TCIMACRO{\TeXButton{EndFrame}{\end{frame}}}%
%BeginExpansion
\end{frame}%
%EndExpansion

%TCIMACRO{\TeXButton{BeginFrame}{\begin{frame}}}%
%BeginExpansion
\begin{frame}%
%EndExpansion

\QTR{frametitle}{Point Estimation}

\begin{itemize}
\item How do we know that a certain choice of estimator $\hat{\theta}_{n}$
is good for estimating $\theta $?

\item One way to look at it is to require that $\hat{\theta}$ is close to $%
\theta $ with respect to some distance (metric).

\item An appropriate distance to choose is the mean squared error (provided
that the estimator $\hat{\theta}_{n}$ has finite variance).
\end{itemize}

\begin{definition}
\textbf{The mean squared error} (MSE) of a statistic $\hat{\theta}_{n}$ is
defined as 
\begin{equation*}
MSE\left( \hat{\theta}_{n},\theta \right) :=\mathbb{E}\left[ \left( \hat{%
\theta}_{n}-\theta \right) ^{2}\right] ,
\end{equation*}%
and more generally $MSE\left( \hat{\theta}_{n},\theta \right) :=\mathbb{E}%
\left[ \left( \hat{\theta}_{n}-\theta \right) \left( \hat{\theta}_{n}-\theta
\right) ^{\prime }\right] $ if $\theta $ is a vector.
\end{definition}

%TCIMACRO{\TeXButton{EndFrame}{\end{frame}}}%
%BeginExpansion
\end{frame}%
%EndExpansion

%TCIMACRO{\TeXButton{BeginFrame}{\begin{frame}}}%
%BeginExpansion
\begin{frame}%
%EndExpansion

\QTR{frametitle}{Point Estimation}

\begin{itemize}
\item We prefer an estimator $\hat{\theta}_{n}$ to a competing estimator $%
\tilde{\theta}_{n}$ iff%
\begin{equation*}
MSE\left( \hat{\theta}_{n},\theta \right) -MSE\left( \tilde{\theta}%
_{n},\theta \right) \;\text{is negative definite matrix.}
\end{equation*}
\end{itemize}

\begin{definition}
\begin{equation*}
MSE\left( \hat{\theta}_{n},\theta \right) =\left( \mathbb{E}\left( \hat{%
\theta}_{n}\right) -\theta \right) \left( \mathbb{E}\left( \hat{\theta}%
_{n}\right) -\theta \right) ^{\prime }+var\left( \hat{\theta}_{n}\right) .
\end{equation*}
\end{definition}

\begin{itemize}
\item Ideally, we would like to minimise $MSE\left( \hat{\theta}_{n},\theta
\right) $ w.r.t. $\hat{\theta}_{n}$ on order to find an optimal estimator.

\item However, joint minimisation is not possible in general.
\end{itemize}

%TCIMACRO{\TeXButton{EndFrame}{\end{frame}}}%
%BeginExpansion
\end{frame}%
%EndExpansion

%TCIMACRO{\TeXButton{BeginFrame}{\begin{frame}}}%
%BeginExpansion
\begin{frame}%
%EndExpansion

\QTR{frametitle}{Point Estimation}

\begin{itemize}
\item Things get a lot simpler if $\hat{\theta}_{n}$ has the additional
property of unbiasedness.
\end{itemize}

\begin{definition}
An estimator $\hat{\theta}_{n}$ of $\theta $ is called \textbf{unbiased} if 
\QTR{Bbb}{E}$\left( \hat{\theta}_{n}\right) =\theta $.
\end{definition}

\begin{itemize}
\item If $\hat{\theta}_{n}$ is unbiased, then $MSE\left( \hat{\theta}%
_{n},\theta \right) =var\left( \hat{\theta}_{n}\right) $.

\item The MSE minimisation criterion implies that: \textit{between two
unbiased estimators we choose the one with the smallest variance}.
\end{itemize}

%TCIMACRO{\TeXButton{EndFrame}{\end{frame}}}%
%BeginExpansion
\end{frame}%
%EndExpansion

%TCIMACRO{\TeXButton{BeginFrame}{\begin{frame}}}%
%BeginExpansion
\begin{frame}%
%EndExpansion

\QTR{frametitle}{Point Estimation}

\begin{itemize}
\item However, one needs to bear in mind that unbiasedness is a compromise:
there exist examples of a biased estimator that has smaller MSE than an
unbiased estimator.

\item Finally, here is another desirable property of estimators in large
samples.
\end{itemize}

\begin{definition}
An estimator $\hat{\theta}_{n}$ of $\theta $ is called \textbf{consistent}
if $\hat{\theta}_{n}\rightarrow _{p}\theta $ as the sample size $n$ tends to 
$\infty $.
\end{definition}

\begin{definition}
$Z_{n}$ converges to $Z$ in probability, written $Z_{n}\rightarrow _{p}Z,$
if for every $\varepsilon >0$ (no matter how small),%
\begin{equation*}
P(|Z_{n}-Z|\geq \varepsilon )\rightarrow 0\text{ as }n\rightarrow \infty .
\end{equation*}
\end{definition}

%TCIMACRO{\TeXButton{EndFrame}{\end{frame}}}%
%BeginExpansion
\end{frame}%
%EndExpansion

%TCIMACRO{\TeXButton{BeginFrame}{\begin{frame}}}%
%BeginExpansion
\begin{frame}%
%EndExpansion

\QTR{frametitle}{Confidence Intervals}

\begin{itemize}
\item We may need more than just a point estimator, we need to know about
the distribution of our estimator, in order to construct Confidence
Intervals and conduct Hypothesis Testing

\item E.g., rather than requiring only an approximation of the value of an
unknown parameter, we might be more interested in constructing (random)
bounds that contain the true unknown parameter with given probability

\item The approach is to construct a random interval that contains the true
parameter with a satisfactory (say 95\%) probability
\end{itemize}

%TCIMACRO{\TeXButton{EndFrame}{\end{frame}}}%
%BeginExpansion
\end{frame}%
%EndExpansion

%TCIMACRO{\TeXButton{BeginFrame}{\begin{frame}}}%
%BeginExpansion
\begin{frame}%
%EndExpansion

\QTR{frametitle}{Confidence Intervals}

\begin{itemize}
\item For the purpose we will need to derive the estimator's exact or large
sample (asymptotic) distribution

\item Many estimators that we will study have asymptotically Normal
distributions

\item Let $X_{1},...,X_{n}$ be a sequence of random variables and let $X$ be
another random variable. Let $F_{n}$ denote the cdf of $X_{n}$ and let $F$
denote the cdf of $X$
\end{itemize}

\begin{definition}
$X_{n}$ converges to $X$ in distribution, written $X_{n}\rightarrow _{d}X$ if%
\begin{equation*}
\lim_{n\rightarrow \infty }F_{n}(x)=F(x)
\end{equation*}%
for all $x$ for which $F$ is continuous.
\end{definition}

%TCIMACRO{\TeXButton{EndFrame}{\end{frame}}}%
%BeginExpansion
\end{frame}%
%EndExpansion

%TCIMACRO{\TeXButton{BeginFrame}{\begin{frame}}}%
%BeginExpansion
\begin{frame}%
%EndExpansion

\QTR{frametitle}{Hypothesis Testing}

\begin{itemize}
\item We might also be interested in testing if the unknown parameters $%
\theta $ satisfy certain suggested restrictions, or not.

\item The restrictions that we declare to be of interest are called the 
\textbf{null hypothesis} (about the parameters $\theta ),$ and this is
usually abbreviated with $H_{0}.$

\item The situation that we imagining to be true when this null hypothesis
is not true is called the \textbf{alternative hypothesis}, and is usually
denoted by $H_{1}.$

\item This may be simply that the null hypothesis is not true, or it may be
more specific.

\item Problem: \textit{Use the data }$X=\left( X_{1},...,X_{n}\right)
^{\prime }$\textit{\ to decide whether}%
\begin{equation*}
H_{0}:\theta \in \Theta _{0}\text{\ \ \ or\ \ \ }H_{1}:\theta \in \Theta _{1}
\end{equation*}%
\textit{where }$\Theta _{0}$\textit{\ and }$\Theta _{1}$\textit{\ are
disjoint subsets of }$\Theta $\textit{.}
\end{itemize}

%TCIMACRO{\TeXButton{EndFrame}{\end{frame}}}%
%BeginExpansion
\end{frame}%
%EndExpansion

\section{The Linear Model}

%TCIMACRO{\TeXButton{BeginFrame}{\begin{frame}}}%
%BeginExpansion
\begin{frame}%
%EndExpansion

\QTR{frametitle}{The linear regression model}

\begin{itemize}
\item The starting point for any econometrics course is the simple linear
regression model

\item Let's first look at a motivating example

\item Suppose we want to value real assets, for example housing price as a
linear function of the property characteristics

\item Suppose we have a sample of $n$ observations for\ the selling price of
a property (in units of \$1,000) and the area of the property in square feet

\item The model is $House$ $price_{i}=\beta $ $Square$ $Feet_{i}+\varepsilon
_{i},$ \ \ $i=1,...,n$

\item Note we are explaining the price with the size of the property and not
the other way around

\item The aim is to estimate the unknown parameter $\beta $
\end{itemize}

%TCIMACRO{\TeXButton{EndFrame}{\end{frame}}}%
%BeginExpansion
\end{frame}%
%EndExpansion

%TCIMACRO{\TeXButton{BeginFrame}{\begin{frame}}}%
%BeginExpansion
\begin{frame}%
%EndExpansion

%TCIMACRO{%
%\FRAME{ftbpF}{4.4693in}{2.8115in}{0pt}{}{}{houseprices.bmp}{%
%\special{language "Scientific Word";type "GRAPHIC";display "USEDEF";valid_file "F";width 4.4693in;height 2.8115in;depth 0pt;original-width 8.0134in;original-height 4.7469in;cropleft "0";croptop "1";cropright "1";cropbottom "0";filename '../Financial Econometrics/houseprices.bmp';file-properties "XNPEU";}}}%
%BeginExpansion
\begin{figure}[ptb]\begin{center}
\includegraphics[
natheight=4.7469in, natwidth=8.0134in, height=2.8115in, width=4.4693in]
{C:/Users/u170828/Dropbox/TEACHING/Econometrics UPF/graphics/houseprices__1.pdf}%
\end{center}\end{figure}%
%EndExpansion

%TCIMACRO{\TeXButton{EndFrame}{\end{frame}}}%
%BeginExpansion
\end{frame}%
%EndExpansion

%TCIMACRO{\TeXButton{BeginFrame}{\begin{frame}}}%
%BeginExpansion
\begin{frame}%
%EndExpansion

\QTR{frametitle}{The Least Squares Estimator}

\begin{itemize}
\item In general, we have a sample of $n$ observations from some dependent
variable $y_{i}$ and some explanatory variable $x_{i}$

\item The model is $y_{i}=\beta x_{i}+\varepsilon _{i},$ $i=1,...,n$

\item The goal: devise a procedure that allows us through observing $y_{i}$
and $x_{i},$ to make a `good' guess for $\beta ,$ by ignoring the random
unobserved component in $\varepsilon _{i}$

\item The least squares estimator for $\beta $ is defined as minimiser of
the sum of squared errors $\hat{\beta}_{n}=\arg \min_{\beta
}\tsum\nolimits_{i=1}^{n}\left( y_{i}-\beta x_{i}\right) ^{2}$

\item Taking first order conditions wrt $\beta $%
\begin{equation*}
\frac{\partial \tsum\nolimits_{i=1}^{n}\left( y_{i}-\beta x_{i}\right) ^{2}}{%
\partial \beta }=-2\tsum\nolimits_{i=1}^{n}\left( y_{i}-\beta x_{i}\right)
x_{i}=0
\end{equation*}

\item The estimator is given by 
\begin{equation*}
\hat{\beta}_{n}=\frac{\tsum\nolimits_{i=1}^{n}y_{i}x_{i}}{%
\tsum\nolimits_{i=1}^{n}x_{i}^{2}}
\end{equation*}
\end{itemize}

%TCIMACRO{\TeXButton{EndFrame}{\end{frame}}}%
%BeginExpansion
\end{frame}%
%EndExpansion

%TCIMACRO{\TeXButton{BeginFrame}{\begin{frame}}}%
%BeginExpansion
\begin{frame}%
%EndExpansion

%TCIMACRO{%
%\FRAME{ftbpF}{2.8755in}{3.0848in}{0pt}{}{}{ssr.png}{%
%\special{language "Scientific Word";type "GRAPHIC";maintain-aspect-ratio TRUE;display "USEDEF";valid_file "F";width 2.8755in;height 3.0848in;depth 0pt;original-width 6.2431in;original-height 6.7014in;cropleft "0";croptop "1";cropright "1";cropbottom "0";filename 'SSR.png';file-properties "XNPEU";}}}%
%BeginExpansion
\begin{figure}[ptb]\begin{center}
\includegraphics[
natheight=6.7014in, natwidth=6.2431in, height=3.0848in, width=2.8755in]
{C:/Users/u170828/Dropbox/TEACHING/Econometrics UPF/graphics/SSR__2.png}%
\end{center}\end{figure}%
%EndExpansion

%TCIMACRO{\TeXButton{EndFrame}{\end{frame}}}%
%BeginExpansion
\end{frame}%
%EndExpansion

%TCIMACRO{\TeXButton{BeginFrame}{\begin{frame}}}%
%BeginExpansion
\begin{frame}%
%EndExpansion

\QTR{frametitle}{The linear regression model}

\begin{itemize}
\item In general, we expect that there is more than one explanatory variable
for the price of a property

\item For instance, the age of the property, the neighbourhood, the number
of bedrooms, the quality of furnishing etc.

$\ \ \ \ \ \ \ \ \ \ y_{i}=\alpha +x_{i1}\beta _{1}+....+x_{ik-1}\beta
_{k-1}+\varepsilon _{i},$ \ $i=1,...,n$

$\qquad \qquad y_{i}=\underset{\mathbf{x}_{i}^{\prime }}{\underbrace{\left[ 
\begin{array}{cccc}
1 & x_{i1} & \cdots & x_{ik-1}%
\end{array}%
\right] }}\underset{\beta }{\underbrace{\left[ 
\begin{array}{c}
\alpha \\ 
\beta _{1} \\ 
\vdots \\ 
\beta _{k-1}%
\end{array}%
\right] }}+\varepsilon _{i},$

where $y_{i}$ is a dependent observed variable, $\mathbf{x}_{i}$ is a $%
k\times 1$ vector of explanatory variables (regressors), $\varepsilon _{i}$
is a r.v., $\alpha $ is an intercept scalar parameter and $\beta $ is a $%
k\times 1$ vector of slope parameters.
\end{itemize}

%TCIMACRO{\TeXButton{EndFrame}{\end{frame}}}%
%BeginExpansion
\end{frame}%
%EndExpansion

%TCIMACRO{\TeXButton{BeginFrame}{\begin{frame}}}%
%BeginExpansion
\begin{frame}%
%EndExpansion

\QTR{frametitle}{The linear regression model}

\begin{itemize}
\item The sum of squared residuals can be written as $\tsum%
\nolimits_{i=1}^{n}\left( y_{i}-\mathbf{x}_{i}^{\prime }\beta \right)
^{2}=\tsum\nolimits_{i=1}^{n}y_{i}^{2}-2\beta ^{\prime
}\tsum\nolimits_{i=1}^{n}\mathbf{x}_{i}y_{i}+\beta ^{\prime
}\tsum\nolimits_{i=1}^{n}\mathbf{x}_{i}\mathbf{x}_{i}^{\prime }\beta $

\item Taking first order conditions wrt the vector $\beta $%
\begin{equation*}
\frac{\partial \tsum\nolimits_{i=1}^{n}\left( y_{i}-\mathbf{x}_{i}^{\prime
}\beta \right) ^{2}}{\partial \beta }=-2\tsum\nolimits_{i=1}^{n}\mathbf{x}%
_{i}y_{i}+2\tsum\nolimits_{i=1}^{n}\mathbf{x}_{i}\mathbf{x}_{i}^{\prime
}\beta
\end{equation*}

\item Solving for the vector $\beta $ yields 
\begin{equation*}
\hat{\beta}_{n}=\left( \tsum\nolimits_{i=1}^{n}\mathbf{x}_{i}\mathbf{x}%
_{i}^{\prime }\right) ^{-1}\tsum\nolimits_{i=1}^{n}\mathbf{x}_{i}y_{i}
\end{equation*}
\end{itemize}

%TCIMACRO{\TeXButton{EndFrame}{\end{frame}}}%
%BeginExpansion
\end{frame}%
%EndExpansion

%TCIMACRO{\TeXButton{BeginFrame}{\begin{frame}}}%
%BeginExpansion
\begin{frame}%
%EndExpansion

%TCIMACRO{%
%\FRAME{ftbpF}{3.1739in}{3.25in}{0pt}{}{}{ssr2d.png}{%
%\special{language "Scientific Word";type "GRAPHIC";maintain-aspect-ratio TRUE;display "USEDEF";valid_file "F";width 3.1739in;height 3.25in;depth 0pt;original-width 6.576in;original-height 6.736in;cropleft "0";croptop "1";cropright "1";cropbottom "0";filename 'SSR2D.png';file-properties "XNPEU";}}}%
%BeginExpansion
\begin{figure}[ptb]\begin{center}
\includegraphics[
natheight=6.736in, natwidth=6.576in, height=3.25in, width=3.1739in]
{C:/Users/u170828/Dropbox/TEACHING/Econometrics UPF/graphics/SSR2D__3.png}%
\end{center}\end{figure}%
%EndExpansion

%TCIMACRO{\TeXButton{EndFrame}{\end{frame}}}%
%BeginExpansion
\end{frame}%
%EndExpansion

%TCIMACRO{\TeXButton{BeginFrame}{\begin{frame}}}%
%BeginExpansion
\begin{frame}%
%EndExpansion

%TCIMACRO{%
%\FRAME{ftbpF}{3.3269in}{2.8885in}{0pt}{}{}{objf.png}{%
%\special{language "Scientific Word";type "GRAPHIC";maintain-aspect-ratio TRUE;display "USEDEF";valid_file "F";width 3.3269in;height 2.8885in;depth 0pt;original-width 7.6181in;original-height 6.6115in;cropleft "0";croptop "1";cropright "1";cropbottom "0";filename 'objf.png';file-properties "XNPEU";}}}%
%BeginExpansion
\begin{figure}[ptb]\begin{center}
\includegraphics[
natheight=6.6115in, natwidth=7.6181in, height=2.8885in, width=3.3269in]
{C:/Users/u170828/Dropbox/TEACHING/Econometrics UPF/graphics/objf__4.png}%
\end{center}\end{figure}%
%EndExpansion

%TCIMACRO{\TeXButton{EndFrame}{\end{frame}}}%
%BeginExpansion
\end{frame}%
%EndExpansion

%TCIMACRO{\TeXButton{BeginFrame}{\begin{frame}}}%
%BeginExpansion
\begin{frame}%
%EndExpansion

\QTR{frametitle}{The linear regression model with matrices}

\begin{itemize}
\item We can equivalently write the model in a matrix form as 
\begin{equation}
y=X\beta +\varepsilon  \label{model}
\end{equation}%
where $y=$ $\left[ 
\begin{array}{c}
y_{1} \\ 
\vdots \\ 
y_{n}%
\end{array}%
\right] $ and $\varepsilon =$ $\left[ 
\begin{array}{c}
\varepsilon _{1} \\ 
\vdots \\ 
\varepsilon _{n}%
\end{array}%
\right] $ are $n\times 1$ vectors, $X=\left[ 
\begin{array}{c}
\mathbf{x}_{1}^{\prime } \\ 
\vdots \\ 
\mathbf{x}_{n}^{\prime }%
\end{array}%
\right] =\left[ 
\begin{array}{cccc}
1 & x_{11} & \cdots & x_{1k-1} \\ 
\vdots & \ddots &  &  \\ 
1 & x_{n1} & \cdots & x_{nk-1}%
\end{array}%
\right] $ is a $n\times k$ matrix, $\beta =\left[ 
\begin{array}{c}
\alpha \\ 
\beta _{1} \\ 
\vdots \\ 
\beta _{k-1}%
\end{array}%
\right] $ is a $k\times 1$ vector.
\end{itemize}

%TCIMACRO{\TeXButton{EndFrame}{\end{frame}}}%
%BeginExpansion
\end{frame}%
%EndExpansion

%TCIMACRO{\TeXButton{BeginFrame}{\begin{frame}}}%
%BeginExpansion
\begin{frame}%
%EndExpansion

\begin{itemize}
\item The OLS estimator $\widehat{\beta }_{n}$ is defined as the minimiser
of following objective function%
\begin{equation*}
\arg \min_{\beta }\left\Vert y-X\beta \right\Vert ^{2}=\arg \min_{\beta
}(y^{\prime }y-2\beta ^{\prime }X^{\prime }y+\beta ^{\prime }X^{\prime
}X\beta )
\end{equation*}

\item Differentiate the objective function wrt the vector $\beta $%
\begin{gather*}
\frac{\partial \left\Vert y-X\beta \right\Vert ^{2}}{\partial \beta }%
=-2X^{\prime }y+2X^{\prime }X\beta =0 \\
\Rightarrow \widehat{\beta }_{n}=(X^{\prime }X)^{-1}X^{\prime }y
\end{gather*}

\item We decompose $y$ into explained and unexplained part: $\ \ \ \ \ \ \ \
\ \ \ \ \ \ \ \ \ \ \ \ \ \ \ \ \ y=\hat{y}+\hat{\varepsilon}=X\widehat{%
\beta }_{n}+\hat{\varepsilon}$

\item We refer to $\hat{\varepsilon}=y-X\widehat{\beta }_{n}$ as 'fitted
residuals'
\end{itemize}

%TCIMACRO{\TeXButton{EndFrame}{\end{frame}}}%
%BeginExpansion
\end{frame}%
%EndExpansion

%TCIMACRO{\TeXButton{BeginFrame}{\begin{frame}}}%
%BeginExpansion
\begin{frame}%
%EndExpansion

\QTR{frametitle}{Housing Price Example}

\begin{itemize}
\item Let's go back to the US\ housing price example

\item We can now include a number of other explanatory variables
(regressors): the number of rooms, the number of bedrooms and the age of the
property

\item Our regression model looks like this

{\small Houseprice}$_{i}=\alpha +\beta _{1}${\small SquareFeet}$_{i}+\beta
_{2}${\small NoRooms}$_{i}+\beta _{3}${\small NoBedrooms}$_{i}+\beta _{4}$%
{\small Age}$_{i}+\varepsilon _{i},$

\item The OLS estimates in this case are

{\small Houseprice}$_{i}=10.37+0.05${\small SquareFeet}$_{i}+6.32${\small %
NoRooms}$_{i}-11.10${\small NoBedrooms}$_{i}+0.43${\small Age}$%
_{i}+\varepsilon _{i},$

\item Interpretation of the coefficients: $\Delta House$ $%
price_{i}=0.05\Delta SquareFeet_{i},$ so that an extra square foot in the
size of the property increases the house of the property by $0.05;$ and
since the price is in units of \$1,000, by $\$50.$
\end{itemize}

%TCIMACRO{\TeXButton{EndFrame}{\end{frame}}}%
%BeginExpansion
\end{frame}%
%EndExpansion

%TCIMACRO{\TeXButton{BeginFrame}{\begin{frame}}}%
%BeginExpansion
\begin{frame}%
%EndExpansion

\QTR{frametitle}{Projection Interpretation}

\begin{itemize}
\item Let's first consider the case when $k=1$ and $n=3.$

\item The Least Squares method minimises the distance between $y$ and $%
X\beta :\min_{\beta }\left\Vert y-X\beta \right\Vert ^{2}$

\item Graphically, the solution to this problem is given by orthogonally
projecting the vector $y$ to the line in $\mathbb{R}^{n}$ defined by $%
c(X)=\{\lambda X:\lambda \in \mathbb{R}\}$

\item Dropping a perpendicular from vector $y$ to the line $c(X)$, we obtain
the orthogonal projection $P_{X}(y)$ of $y.$

\item Since $P_{X}(y)$ is the closest point of $c(X)$ to $y,$ $P_{X}(y)=X%
\widehat{\beta }$

\item Also, since $y-P_{X}(y)\perp X\Rightarrow X^{\prime }(y-P_{X}(y))=0$

\item So, $X^{\prime }y=X^{\prime }P_{X}(y),$ and substituting the above
expression for $P_{X}(y),$ we obtain $X^{\prime }y=X^{\prime }X\widehat{%
\beta }\Rightarrow \widehat{\beta }=(X^{\prime }X)^{-1}X^{\prime }y.$
\end{itemize}

%TCIMACRO{\TeXButton{EndFrame}{\end{frame}}}%
%BeginExpansion
\end{frame}%
%EndExpansion

%TCIMACRO{\TeXButton{BeginFrame}{\begin{frame}}}%
%BeginExpansion
\begin{frame}%
%EndExpansion

\QTR{frametitle}{Graphical Representation}

%TCIMACRO{%
%\FRAME{ftbpF}{4.1727in}{3.0286in}{0pt}{}{}{orthogonal-page0001.bmp}{%
%\special{language "Scientific Word";type "GRAPHIC";display "USEDEF";valid_file "F";width 4.1727in;height 3.0286in;depth 0pt;original-width 5.9733in;original-height 5.6196in;cropleft "0";croptop "1";cropright "1";cropbottom "0";filename '../Financial Econometrics/orthogonal-page0001.bmp';file-properties "XNPEU";}}}%
%BeginExpansion
\begin{figure}[ptb]\begin{center}
\includegraphics[
natheight=5.6196in, natwidth=5.9733in, height=3.0286in, width=4.1727in]
{C:/Users/u170828/Dropbox/TEACHING/Econometrics UPF/graphics/orthogonal-page0001__5.pdf}%
\end{center}\end{figure}%
%EndExpansion

%TCIMACRO{\TeXButton{EndFrame}{\end{frame}}}%
%BeginExpansion
\end{frame}%
%EndExpansion

%TCIMACRO{\TeXButton{BeginFrame}{\begin{frame}}}%
%BeginExpansion
\begin{frame}%
%EndExpansion

\QTR{frametitle}{Projection Interpretation}

\begin{itemize}
\item $P_{X}(y)=X\widehat{\beta }=X(X^{\prime }X)^{-1}X^{\prime
}y\Rightarrow P_{X}=X(X^{\prime }X)^{-1}X^{\prime }.$

\item $\widehat{\beta }$ is the coordinate of $y$ in the direction of $X$
(i.e. on the line $c(X)$).

\item The coordinate of $y$ is given my projecting $y$ on the orthogonal
complement of $c(X),$ given by $c(X)^{\perp }=\{z\in \mathbb{R}%
^{n}:z^{\prime }X=0\}.$

\item Projecting $y$ on $c(X)^{\perp }$ gives a projection $%
y-P_{X}(y)=(I-X(X^{\prime }X)^{-1}X^{\prime })y=P_{X^{\perp }}(y).$

\item $y=P_{X}(y)+P_{X^{\perp }}(y)=X\widehat{\beta }+\widehat{\varepsilon }=%
\hat{y}+\hat{\varepsilon}$

\item We would refer to the projection matrix $P_{X^{\perp }}$ as $M_{X}$

\item In the general case, $X$ is a $n\times k$ matrix with rank $k,$ $c(X)$
is a $k-$dimensional subspace of $\mathbb{R}^{n};$ the projection matrices $%
P_{X}$ and $M_{X}$ have the same algebraic form as in the case $k=1;$ only
now they project on more complicated subspaces.
\end{itemize}

%TCIMACRO{\TeXButton{EndFrame}{\end{frame}}}%
%BeginExpansion
\end{frame}%
%EndExpansion

%TCIMACRO{\TeXButton{BeginFrame}{\begin{frame}}}%
%BeginExpansion
\begin{frame}%
%EndExpansion

\QTR{frametitle}{Multiple Regression Graphical Representation}

%TCIMACRO{%
%\FRAME{ftbpF}{4.625in}{2.8089in}{0pt}{}{}{orthogonal2-page0001.bmp}{%
%\special{language "Scientific Word";type "GRAPHIC";display "USEDEF";valid_file "F";width 4.625in;height 2.8089in;depth 0pt;original-width 6.3062in;original-height 5.4474in;cropleft "0";croptop "1";cropright "1";cropbottom "0";filename '../Financial Econometrics/orthogonal2-page0001.bmp';file-properties "XNPEU";}}}%
%BeginExpansion
\begin{figure}[ptb]\begin{center}
\includegraphics[
natheight=5.4474in, natwidth=6.3062in, height=2.8089in, width=4.625in]
{C:/Users/u170828/Dropbox/TEACHING/Econometrics UPF/graphics/orthogonal2-page0001__6.pdf}%
\end{center}\end{figure}%
%EndExpansion

%TCIMACRO{\TeXButton{EndFrame}{\end{frame}}}%
%BeginExpansion
\end{frame}%
%EndExpansion

%TCIMACRO{\TeXButton{BeginFrame}{\begin{frame}}}%
%BeginExpansion
\begin{frame}%
%EndExpansion

\QTR{frametitle}{Reading}

Relevant chapters from the textbook: Chapter 1-3.

%TCIMACRO{\TeXButton{EndFrame}{\end{frame}}}%
%BeginExpansion
\end{frame}%
%EndExpansion

\end{document}
